% ============================================================================
% CHAPTER 6 - CONCLUSION   (template file chapter_9.tex)
% ============================================================================

\section{Summary of Findings}

This thesis set out to turn a recorded Bengali-English code-mixed whiteboard
lecture into a lecture note that a student can use. The system was built as four
stages and each stage was measured separately. The findings are summarised below in
the order they were produced.

\textbf{Off-the-shelf speech recognition does not work on this speech, and the
failure is specific.} On six held-out lectures, Whisper large-v3-turbo has a median
character error rate of 67.7 per cent and a median word error rate of 93.9 per cent.
It translates the speech into English instead of transcribing it, and it falls into
repetition loops on thirteen of the 177 test clips.

\textbf{A small adapter on the same model fixes most of that.} Trained on 4.10 hours
of hand-checked Banglish from 22 lectures, LoRA adapters bring the character error
rate down to 15.8 and 16.0 per cent in two random seeds, and the word error rate to
42.5 and 41.7 per cent. The adapted model is better on 164 and 167 of 177 clips with
p below 1e-30, and it needs no decoding safeguard, since it loops on one clip in 177.

\textbf{The gain also holds for a lecturer the model has never heard, at a lower
level.} Holding out each lecturer in turn and training on the other two cuts the
character error rate from 72.8 to 50.3 per cent, significant in all six runs. This
is the harder question and the worse number, and both designs are reported.

\textbf{Training at the selected settings is not stable.} At the learning rate the
pre-registered tuning procedure chose, one of two seeds collapsed into repetition and
produced output worse than doing nothing. The stable configuration reported beside it
is the runner-up rate from the same table, and it was run after the divergence was
seen.

\textbf{The whiteboard can be reconstructed from the video without inventing
anything.} Tiling the frame and taking each tile from a moment when the lecturer was
not in front of it gives boards in which every pixel is unmodified camera output. The
median board has 97.7 per cent of its tiles fully clear. A human check of 98 boards
found 82 complete, with writing genuinely lost on 9 and the remaining 7 explained by
the source video rather than by the method.

\textbf{Reading the board is a question of prompting, not of model size.} The same
vision-language model on the same images goes from 31.2 to 88.8 per cent recall of
hand-verified board content when the prompt asks for a full transcription instead of
keywords, and from none of 67 numbers to 65 of 67. A model less than half the size
reaches 85.4 per cent with the same prompt. The prompt is worth 57.6 points and the
model size 5.3.

\textbf{Board content can be made to reach the notes.} The annotated notes carry 89.1
per cent of the board items against 37.2 per cent for a conventional summarisation
prompt, and they match a variant that simply pastes the extracted board text, while
remaining a readable document with numbered regions, verified quotations and a
step-by-step explanation.

\textbf{Several things did not work, and that is part of the result.} Dual recogniser
fusion gives 0.7 percentage points at p = 0.32. Biasing the decoder towards board
words halves term recall at every non-zero strength. Using occlusion as a pointer to
align speech with board regions is not supported, at p = 0.054 and p = 0.159. The
display clean-up that makes boards look better makes the model read them very slightly
worse. Fine-tuning the speech model does not measurably improve board-content recall
in the notes.

\textbf{Measures built out of English strings mislead in this setting.} This appeared
twice, independently. A term-based F1 over a fixed English word list swings 8.1 points
between two training runs on identical data, and it moves against transcription
quality because a better Banglish transcript contains fewer English words. In the note
evaluation, the transcript with 67.7 per cent character error contains significantly
more of the English answer key items than the one with 15.8 per cent, for the same
mechanical reason.

\section{Contributions to the Field}

The contributions of this thesis, stated as narrowly as the evidence allows, are the
following.

\textbf{A hand-checked corpus of code-mixed classroom speech in romanized Banglish.}
Forty-four lecture recordings totalling 8.33 hours, of which 28 lectures and 5.15
hours are transcribed word by word against the audio according to a written
transcription standard, with 609 timed segments and about 39,100 words. As far as
could be found, no comparable corpus of spontaneous Bengali-English classroom speech
transcribed in the Roman alphabet existed. The nearest published corpora in this
language pair are synthetic.

\textbf{A measured speech recognition result for this variety, under a strict
evaluation design.} Whole lectures are held out, references are human, the settings
come from a procedure written down before any run, and two seeds are reported for
every configuration including the one that diverged. Two designs are reported and
labelled, one with the test lecturers heard in training and one with an entirely
unseen lecturer.

\textbf{A whiteboard reconstruction that generates nothing, read into text by a
vision-language model and scored against hand-verified answer keys.} The reconstruction
itself is engineering assembled from known parts, and the thesis says so. What is new
is the combination with a vision-language reader and the measurement of it, including
45 boards and 436 items verified by hand and 98 boards checked for missing writing.

\textbf{A demonstration that the prompt, not the model size, is what makes a
vision-language model read a handwritten board.} 57.6 points against 5.3 points, with
the choice re-checked on a held-out half of the boards.

\textbf{A note generator whose claims can be checked.} Board content is supplied as
text rather than recalled, non-lecture material is confined to a labelled box,
quotations are verified word for word against the transcript and deleted when they
fail, and references to numbered regions of the board are checked against the regions
that exist. Each of these is counted in the output file of every page.

\textbf{A methodological finding about evaluating code-mixed systems.} Measures
assembled from English strings reward a model that translates and penalise one that
transcribes. This was observed twice in unrelated measures within one project, it is
consistent with the published literature on code-switched evaluation, and it is
reported here with the mechanism and the numbers rather than as an aside.

\textbf{A set of measured negative results.} Fusion, decoder bias, occlusion-as-pointer
and the display clean-up were all built, all measured and all reported as they came
out.

\section{Recommendations for Future Work}

\subsection{The work that would most improve the thesis}

\textbf{A reader study.} This is the single most valuable missing piece. Nothing in
this thesis measures whether the notes are good to read or whether what they say is
true, and a factual error in the generated prose has already been found by reading.
The study should show about twenty students the old notes and the new notes for the
same lecture and ask which helps more, together with ratings for usefulness,
correctness and readability, giving a paired before-and-after result rather than a
score. A second arm should compare notes built from the off-the-shelf transcript with
notes built from the fine-tuned one, which is the only way to test the part of the
speech improvement that board-content recall cannot see.

\textbf{A benchmark against a neighbouring model.} Running a public Bengali-English
code-switched model on our own test clips would turn the discussion in Section 5.5.1
into a table. It costs about twenty minutes on any suitable machine and it was not
done before submission.

\textbf{Measuring post-editing effort.} The economic argument in Section 3.8.3 rests on
a projection. Timing transcribers as they correct real model drafts would replace it
with a measurement, and it is the number a department would actually want.

\subsection{Work on the speech side}

\textbf{More speech, especially from under-represented lecturers.} The lecturer with
0.7 hours in training has a test lecture at 52 per cent character error while the
others reach 12 to 16. Collecting more from that lecturer would test whether the cause
is the data or the speaker.

\textbf{A scaling curve on the current corpus.} Planned and cancelled when the corpus
closed. It would answer how much more transcription is worth buying.

\textbf{Stabilising the training.} The divergence at the selected learning rate is
reproducible enough to study. Warmup, gradient clipping or a lower alpha are the
obvious things to try, and a configuration that never diverges would be worth more than
a small gain in error rate.

\textbf{A spelling-tolerant metric as standard.} The variants computed in Section
\ref{sec:asr-spelling} were used only to answer a question. Reporting them alongside
plain error rates, and applying the same idea to the note metric, would reduce the
measurement problem described throughout this thesis.

\subsection{Work on the vision and note side}

\textbf{Hold the era open until a wipe finishes.} Four of the 98 boards were cut in the
middle of an erase. This is a small change with a known cause.

\textbf{Ignore board furniture.} A sticker that appears unchanged in every era of a
lecture should not be given a numbered box.

\textbf{Evaluate the box finding.} No layout ground truth exists, so the step that
produces the numbered regions is the one part of the pipeline with no number attached
to it.

\textbf{Notes in Bengali script.} Dropped from this thesis deliberately, because it
would be produced by translation alone and would add checking work without adding a
contribution. It is straightforward to add if a reader wants it.

\textbf{A larger notes model.} A 32 billion parameter model was intended for comparison
and was dropped for disk space. It would only change the writing, not the facts, since
it reads the extracted board text and never the board.

\subsection{Closing remark}

The system described here does what it was built to do, and the report says exactly how
well. A student can put in a recorded Banglish whiteboard lecture and get out a note
containing the lecturer's words and the lecturer's board, with the board drawn from
real pixels and the quotations checked against the transcript. The speech model is good
enough to be a draft and not good enough to be trusted unchecked. The board reading is
strong and the reason it is strong is the prompt. Whether the notes are genuinely
useful to a student is the one question this thesis cannot answer, and it is the first
thing that should be measured next.
